% =============================================================================
\section{Complete Proofs}\label{app:proofs}
% =============================================================================

This appendix collects complete proofs for all theoretical results stated in the main text.

\subsection{Proof of \texorpdfstring{\cref{prop:mono_beta}}{Proposition: Monotonicity}}

\begin{proof}
Let $0<\beta_1\le \beta_2\le 1$. Define $A(\beta) := \{\ell:\cI_{\le \ell}(r)\ge \beta\,\cI_{\mathrm{tot}}(r)\}$.
Because $\beta_1\le \beta_2$, for any $\ell$, $\cI_{\le \ell}(r)\ge \beta_2 \cI_{\mathrm{tot}}(r)$ implies $\cI_{\le \ell}(r)\ge \beta_1 \cI_{\mathrm{tot}}(r)$. Hence $A(\beta_2)\subseteq A(\beta_1)$. Taking minima over nested sets: $\min A(\beta_1)\le \min A(\beta_2)$, i.e., $\ECL_{\beta_1}(r)\le \ECL_{\beta_2}(r)$.
\end{proof}

\subsection{Proof of \texorpdfstring{\cref{thm:exp_decay}}{Theorem: Exponential Decay Bound}}

\begin{proof}
Under the assumption $\cI(d;r) \le C e^{-\lambda d}$ with multiplicity $n_d \le 2$ for $d \ge 1$ and $n_0 = 1$:
\[
\cI_{\mathrm{tot}}(r) \le C + 2C \sum_{d=1}^{L} e^{-\lambda d} = C + \frac{2C e^{-\lambda}}{1-e^{-\lambda}} \le C + \frac{2C}{\lambda}
\]
using the geometric series bound $\sum_{d=1}^\infty e^{-\lambda d} = e^{-\lambda}/(1-e^{-\lambda}) \le 1/\lambda$ for $\lambda > 0$.

The tail influence beyond radius $\ell$ is bounded by:
\[
\sum_{d=\ell+1}^{L} n_d \cI(d;r) \le 2C \sum_{d=\ell+1}^{\infty} e^{-\lambda d} = \frac{2C e^{-\lambda(\ell+1)}}{1-e^{-\lambda}} \le \frac{2C}{\lambda} e^{-\lambda \ell}.
\]

ECL$_\beta(r)$ is the smallest $\ell$ such that $\cI_{\le \ell}(r) \ge \beta \cI_{\mathrm{tot}}(r)$, equivalently such that $\cI_{>\ell}(r) \le (1-\beta)\cI_{\mathrm{tot}}(r)$. A sufficient condition is:
\[
\frac{2C}{\lambda} e^{-\lambda \ell} \le (1-\beta)\cI_{\mathrm{tot}}(r).
\]
Solving for $\ell$:
\[
\ell \ge \frac{1}{\lambda}\log\frac{2C}{\lambda(1-\beta)\cI_{\mathrm{tot}}(r)} = \frac{1}{\lambda}\log\frac{1}{1-\beta} + \frac{1}{\lambda}\log\frac{2C}{\lambda \cI_{\mathrm{tot}}(r)}.
\]
Since $\ECL_\beta(r)$ is the minimum integer $\ell$ satisfying this, adding 1 for the ceiling gives the stated bound.
\end{proof}

\subsection{Proof of \texorpdfstring{\cref{prop:ecl_not_blackwell}}{Proposition: ECL ordering does not imply Blackwell dominance}}

\begin{proof}[Counterexample]
Consider $L=3$, $r=2$, and two models $f,g:\{0,1\}^3 \to \bbR^2$. Define:
\begin{align*}
f(x) &= (x_1 + x_3,\; x_2), \\
g(x) &= (x_2,\; x_1 \oplus x_3),
\end{align*}
where $\oplus$ denotes XOR. Under uniform $P_X$ and single-coordinate substitution perturbation:

For $f$: $\cI(1;2) = \cI(3;2) = 1$, $\cI(2;2) = 1$, so $\ECL_{0.5}(f,2) = 0$, $\ECL_{1.0}(f,2) = 1$.

For $g$: $\cI(1;2) = \cI(3;2) = 1$, $\cI(2;2) = 1$, so the ECL profiles are identical.

Yet neither model Blackwell-dominates the other. Take the target $Y := x_1 + x_3 \pmod 4$. Then $Y$ has four distinct values and is recoverable from $f(X)$ (the first coordinate of $f$ is exactly $x_1 + x_3 \in \{0,1,2\}$), but only the parity of $Y$ is recoverable from $g(X)$ (the second coordinate of $g$ is $x_1 \oplus x_3$, which collapses $\{(0,1),(1,0)\}$ together). Hence there is no Markov kernel $K$ with $g(X) \stackrel{d}{=} K(f(X))$ that simultaneously preserves $P(Y \mid g(X))$, ruling out Blackwell domination of $g$ by $f$. Symmetrically, take $Y' := x_1 \oplus x_3$; then $Y'$ is recoverable from $g(X)$ but $f(X)$ retains only the joint sum $x_1 + x_3$, whose parity is $Y'$ but whose value carries strictly more information than $Y'$ alone, so $f(X)$ also fails to be a deterministic function of $g(X)$, ruling out Blackwell domination of $f$ by $g$. ECL profiles are therefore equal while neither model dominates the other in the Blackwell sense.
\end{proof}

\subsection{Proof of \texorpdfstring{\cref{thm:minimax}}{Theorem: Sample-Complexity Lower Bound}}

\begin{proof}
\textbf{Lower bound.} Consider two models in $\cF_M$: $f_0$ with constant influence $\cI(i;r) = c$ for all $i$, and $f_1$ with $\cI(i;r) = c + \epsilon$ for $|i-r| \le \ell_0$ and $\cI(i;r) = c - \epsilon'$ otherwise, where $\epsilon, \epsilon'$ are chosen so that $\ECL_\beta(f_0) = \ell_0$ and $\ECL_\beta(f_1) = \ell_0 + 1$ with margin $\gamma = \Theta(\epsilon)$. Each pair of models induces a pair of observation distributions over the $n$ i.i.d.\ sequence-perturbation samples whose KL divergence is $O(n\epsilon^2/M^2)$ since the influence variables are bounded in $[0,M]$. By Le Cam's two-point method, any estimator $\hat{\ell}$ achieving worst-case error probability at most $1/3$ for both $f_0, f_1$ requires $\mathrm{KL}(P_0^{\otimes n} \| P_1^{\otimes n}) \ge c'$ for some absolute constant $c'$, which forces $n = \Omega(M^2/\gamma^2)$. This yields the stated sample-complexity lower bound.

\textbf{Matching upper bound.} By \cref{thm:ecl_stability}, $\widehat{\ECL}_\beta = \ECL_\beta$ when $\max_i |\widehat{\cI}(i;r) - \cI(i;r)| \le \gamma/(4L)$. Applying \cref{thm:bernstein} (Bernstein's inequality for bounded variables) and a union bound over the $L$ positions, this holds with probability $\ge 1-\delta$ when $n \ge C\,L^2 M^2 \gamma^{-2}\,\log(L/\delta)$ for an absolute constant $C$, matching the lower bound up to logarithmic factors in $L$ and $1/\delta$.
\end{proof}

\subsection{Proof of \texorpdfstring{\cref{prop:composition}}{Proposition: Compositional Bound}}

\begin{proof}
Let $h:\cX \to \cH$ and $g:\cH \to \cZ$ with $f = g \circ h$. The output $f(x)_r = g(h(x))_r$ depends on intermediate positions $\{j : |j - r| \le R_g\}$ in $\cH$. Each intermediate position $j$ depends on input positions $\{i : |i - j| \le R_h\}$ in $\cX$. By the triangle inequality, input position $i$ can influence output position $r$ only if there exists intermediate $j$ with $|j-r| \le R_g$ and $|i-j| \le R_h$, implying $|i-r| \le R_h + R_g$. Hence $\cI(i;r) = 0$ whenever $|i-r| > R_h + R_g$, and \cref{prop:ecl_local} gives $\ECL_\beta(r) \le R_h + R_g$.

For a stack of $D$ layers with kernel sizes $k_l$ and strides $s_l$ (all stride one for the dilation case), the per-layer receptive fields compose additively after taking strides into account: $R^* = 1 + \sum_{l=1}^D (k_l - 1) \prod_{m<l} s_m$. For dilated convolutions with constant kernel size $k$ and dilation rates $d_1, \dots, d_D$ (all strides equal to one), this reduces to $R^* = 1 + \sum_{l=1}^D (k-1) d_l$.
\end{proof}
